\textbf{Время:} резерв

\textbf{Тезис:} продольный радиус и lambda относятся к конкретной процедуре анализа.

\textbf{Визуал:} отдельно назвать, какая панель взята из 3D-аппроксимации, а какая из 1D.

\textbf{Текст:} R-long характеризует продольную структуру выбранной области испускания в принятой системе отсчёта. На него влияют отбор и динамика; это не исходная длина ядра.

Lambda справа относится к одномерной аппроксимации. Этот параметр нельзя автоматически приравнивать силе сигнала другого фита. Для экспериментальной выборки важны примеси и разрешение. В чистом модельном наборе такие эффекты не возникают сами собой: нужно отдельно исследовать распады, форму источника и диапазон аппроксимации.

Ошибки оставлены исходными. Без событий и установленного алгоритма неопределённостей не будем давать новую численную оценку значимости архивных трендов.
